\begin{document}

\section{时空组合性的编程范式}\label{a-programming-paradigm-for-spatiotemporal-composability}

Yifan Shi1,2, Wei Zhang1, Tianyi Cui2

1Peking University 2DeepSeek-AI

\subsection{摘要}\label{abstract}

现代软件——从插件系统到自演化智能体框架——日益依赖于动态组合，但其形
式化基础仍不成熟。我们识别出该问题的两个正交维度：时间维组合性，即在
组件被移除时彻底撤回该组件所产生的副作用的能力；以及空间维组合性，即
声明并以响应式方式管理组件间依赖的能力。我们通过将经典的效应与余效应
概念提升为运行时机制来应对这两个维度。具体而言，我们形式化了可逆效
应：每一次上下文变换都带有由运行时追踪的逆变换。我们形式化了响应式余
效应：上下文的每一次变化都会依据其余效应规格通知组件。我们将效应上下
文与余效应上下文统一为单一的上下文类型，由此构成一个编程范式。随后，
我们把这些机制融合为组件的概念，并给出动态组合的演算；该演算的元理论
将时空组合性从单个组件推广到由相互交织的组件构成的整个系统。我们在
Cordis 中实现了这些思想——Cordis 是一个时空组合性的元框架，它提供一
个带有效应追踪与余效应解析的核心库，以及一个支持配置对账与热模块替换
的声明式组件加载器。

\subsection{目录}\label{contents}

\begin{enumerate}
\def\labelenumi{\arabic{enumi}.}
\tightlist
\item
  引言 . . . . . . . . . . . . . . . . . . . . . . . . . . . . .
  . . . . . . . . . . . 4\\
  1.1. 组合性的维度 . . . . . . . . . . . . . . . . . . .
  . . . . . . . . . . . . . . . . . . . . .\\
  1.2. 动机示例 . . . . . . . . . . . . . . . . . . . . . . .
  . . . . . . . . . . . . . . . . . . . . . . . . . . . . . . . . . . .
  . . 4\\
  1.2.1. 插件系统 . . . . . . . . . . . . . . . . . . . . . . . .
  . . . . . . . . . . . . . . . . . . . . . . . . . . . . . . . . . .
  4\\
  1.2.2. 自演化智能体框架 . . . . . . . . . . . . . . . . .
  . . . . . . . . . . . . . . . . . . . 5\\
  1.2.3. 粗粒度的变通方案 . . . . . . . . . . . . . . . . .
  . . . . . . . . . . . . . . . . . . . . . . . 5\\
  1.3. 贡献 . . . . . . . . . . . . . . 6\\
\item
  预备知识 . . . . . . . . . . . . . . . . . . . . . . . . . . . .
  . . . . . . . . . . . . . . . . . . . . . . . . . . . . . . . . . . .
  . . . . . . . . . . . . . . . . . . . 7\\
  2.1. 效应 . . . . . . . . . . . . . . . . . . . . . . . . . . . . .
  . . . . . . . . . . . . . . . . . . . . . . . . . . . . . . 7\\
  2.2. 余效应 . . . . . . . . . . . . . . 7\\
  2.3. 与动态组合的关系 . . . . . . . . . . . . . .
  . . . . . . . . . . . . . . . . . . . . . . . . . . . 8\\
\item
  可逆效应与响应式余效应 . . . . . . . . . . . . . . .
  . . . . . . . . . . . . . . . . . . . . . . . . . . . . . . . . . . .
  . 9\\
  3.1. 可逆效应 . . . . . . . . 9\\
  3.1.1. 效应上下文 . . . . . . . . . . . . . 9\\
  3.1.2. 可逆效应函数 . . . . . . . . . . . . . . . . . .
  . . . . . . . . . . . . . . . . . . . . . . . . . . . . . . . 12\\
  3.1.3. 效应的独立性 . . . . . . . . . . . . . . . . . . . .
  . . . . . . . . . . . . . . . . . . . . . . . . . . . . . . . . . . .
  . . 15\\
  3.2. 响应式余效应 . . . . . . . . . . . . . . . . . . . . . . . .
  . . . . . . . . . . . . . . . . . . . . . . . . . . . . . . . . . . .
  . . . . . . . . . . . . 17\\
  3.2.1. 余效应上下文 \ldots.. . . . . . . . . . . . . . . . . . . .
  . . . . . . . . . . . . . . . . . . . . . 18\\
  3.2.2. 规格与通知 . . . . . . . . . . . . . . . .
  . . . . . . . . . . . . . . . . . . . . . . . . . . . . . 19\\
  3.2.3. 隔离与拦截 . . . . . . . . . . . . . . . . . .
  . . . . . . . . . . . . . . . . . . . . . . . . . . . . . . . 20\\
  3.3. 上下文范式 . . . . . . . . . . . . . . . . . . . . . .
  . . . . . . . . . . . . . . . . . . . . . . . . . . . . . . . . . . .
  . . . . . 22\\
  3.3.1. 统一上下文 . . . . . . . . . . . . . . . . . . . . . . . .
  . . . . . . . . . . . . . . . . . . . 22\\
  3.3.2. 观测等价性 . . . . . . . . . . . . . . . . . . .
  . . . . . . . . . . . . . . . . . . . . . . . . . . . . 23\\
  3.3.3. 上下文范式的定位 . . . . . . . . . . . . . . . .
  . . . . . . . . . . . . . . . . 27\\
\item
  动态组合的演算 . . . . . . . . . . . . . . . . . .
  . . . . . . . . . . . . . . . . . . . . . . . . . . . . . . . . . . .
  . 28\\
  4.1. 组件与纤维 . . . . . . . . . . . . . . . . . . . . . .
  . . . . . . . . . . . . . . . . . . . . . . . . . . . . . . . . . . .
  . 28\\
  4.2. 基础演算 . . . . . . . . 30\\
  4.3. 进行中的转移 . . . . . . . . . . . . . . . . . . . . .
  . . . . . . . . . . . . . . . . . . . . . . . . 33\\
  4.3.1. 撤回 . . . . . . . . . . . . . . . . . . . . . . . . . .
  . . . . . . . . . . . . . . . . . . . . . . . . 34\\
  4.3.2. 迭代 . . . . . . . 35\\
  4.3.3. 异步 . . . . . . . . . . 37\\
  4.3.4. 失败 . . . . . . . 37\\
  4.4. 元理论 . . . . 38\\
  4.4.1. 保持 . . . . . . . 42\\
  4.4.2. 时间维组合性 . . . . . . . . . . . . . . . . . . . .
  . . . . . . . . . . . . . . . . . . . . 43\\
  4.4.3. 空间维组合性 . . . . . . . . . . . . . . . . . . . . .
  . . . . . . . . . . . . . . . . . . . . . . . . . . 45\\
  4.4.4. 进展 . . . . . 47\\
  4.4.5. 汇合 . . . . . . . . . . . . . . 49\\
\item
  实现与案例研究 . . . . . . . . . . . . . . . . . . . .
  . . . . . . . . . . . . . . . . . . . . . . . . . . . . . . . . . .
  54\\
  5.1. 核心库 . . . . . . 54\\
  5.1.1. 效应追踪 . . . . . . . . . . . . . . . . . . . . . . . .
  . . . . . . . . 56\\
  5.1.2. 余效应操作 . . . . . . 57\\
  5.1.3. 组件生命周期 . . . . . . . . . . . . . . . . . . . . . .
  . . . . . . . . . . . . . . . . . . . . . 58\\
  5.1.4. 上下文访问 . . . . . . . 61\\
  5.2. 组件加载器 . . . . . . 61\\
  5.2.1. 声明式配置 . . . . . . . . . . . . . . . . . . .
  . . . . . . . . . . . . . . . . . . . 62\\
  5.2.2. 热模块替换 . . . . . . . . . . . . . . . . . . . .
  . . . . . . . . . . . . . . . . . . 64\\
  5.3. 案例研究：Koishi . . . . . . 66\\
\item
  讨论 . . . . . . . . . . . . . . . . . . . . . . . . . . . . . .
  . . . . . . . . . . . . . . . . . . . . . . . . . . . . . . . . . . .
  . . 67\\
  6.1. 系统边界 . . . . . . 67\\
  6.2. 服务多路复用 . . . . . . . . . . . . . 68\\
  6.3. 访问控制与沙箱 . . . . . . . . . . . . . . . . . .
  . . . . . . . . . . 69\\
  6.4. 语言无关性与选择 . . . . . . . . . . . . . . .
  70\\
  6.5. 相互依赖与组件粒度 . . . . . . . . . .
  . . . . . . . 71\\
  6.6. 依赖类型化与版本化 . . . . . . . . . 72\\
  6.7. 与语言和操作系统的协同设计 . . . . . . . . .
  . . . . . . . . . . . . . . . . . . . . . . . 73\\
\item
  相关工作 . . . . . . . . . . . . . . 74\\
  7.1. 效应与余效应系统 . . . . . . 74\\
  7.2. 编程范式 . . . . . . . . . 75\\
  7.3. 时间维组合性 . . . . . . . . 76\\
  7.4. 空间维组合性 . . . . . . 78\\
\item
  结论 . . . . . . . . 79\\
  参考文献 . . . . . . . . . . . 80
\end{enumerate}

